% !TeX root = ../../Book3.tex
% 纲要标题：### 24.3 Proj 构造
% 纲要位置：计划纲要.md，第 2263 行。

\section{Proj 构造}
\label{b3:sec:2403}


% 纲要要点（供目录核对）：
%
% * 齐次素理想
% * 标准开集与射影空间
% * 分次模和射影几何的联系
%

\subsection{齐次素理想}
\label{b3:subsec:2403:01}

若方程齐次，缩放坐标不会改变其零点条件。直接保留所有仿射点会把同一方向重复记录，并留下不适合表示方向的原点。射影构造因此选取齐次素理想，同时排除包含全部正次数元素的那些素理想。

\begin{definition}\label{b3:def:proj}
设 \(S=\bigoplus_{n\ge0}S_n\) 为交换正分次环，\(S_+=\bigoplus_{n>0}S_n\)。定义
\[
\operatorname{Proj}S
=\{\,\mathfrak p\subset S:\mathfrak p\text{ 为齐次素理想且 }
S_+\not\subseteq\mathfrak p\,\}.
\]
对齐次理想 \(I\)，令 \(V_+(I)=\{\mathfrak p:I\subseteq\mathfrak p\}\)，以这些集合为闭集。对正次数齐次元 \(f\)，记 \(D_+(f)=\operatorname{Proj}S\setminus V_+((f))\)。
\end{definition}

零理想与单位理想分别给全空间和空集；理想和及乘积分别对应闭集的任意交和有限并，证明与素谱相同。各 \(D_+(f)\) 构成开集基，且 \(D_+(f)\cap D_+(g)=D_+(fg)\)。排除 \(S_+\) 的条件恰好保证每个点至少属于其中一个基本开集。

\ProseParagraphBreak
在 \(k[x_0,\ldots,x_r]\) 的标准分次下，非零向量 \(a=(a_0,\ldots,a_r)\in k^{r+1}\) 给出齐次素理想
\[
\ker\bigl(k[x_0,\ldots,x_r]\longrightarrow k[t],\
x_i\longmapsto a_it\bigr).
\]
同比例向量给相同理想；选择 \(a_i\ne0\) 后，坐标比 \(a_j/a_i\) 又确定该点。这恢复通常的 \(k\) 值射影点。Proj 还包含非闭点及可能具有更大剩余域的闭点，不能只用 \(k\) 中向量穷尽一般域上的全部点。

\subsection{标准开集与射影空间}
\label{b3:subsec:2403:02}

\begin{theorem}\label{b3:thm:proj-affine-charts}
设 \(S\) 为交换正分次环，\(f\in S_d\)、\(d>0\) 为齐次元，记 \(S_{(f)}=(S_f)_0\)。则
\[
D_+(f)\cong\operatorname{Spec}S_{(f)}.
\]
这些仿射结构在交集上相容，因而唯一确定 \(\operatorname{Proj}S\) 的概形结构。这里 \(S_{(f)}\) 由分式 \(a/f^m\)、\(a\in S_{md}\) 组成，绝不能用整个 \(S_f\) 代替它。
\end{theorem}
\begin{proof}
从 \(\mathfrak p\in D_+(f)\) 取 \((\mathfrak pS_f)_0\) 得 \(S_{(f)}\) 的素理想。说明逆向的存在唯一性时，令 \(B=(S_f)_0\)，取 \(\mathfrak q\in\operatorname{Spec}B\)，考察分次环
\[
C=S_f\otimes_B\kappa(\mathfrak q).
\]
它的零次部分为 \(\kappa(\mathfrak q)\)，且 \(f\) 仍是次数 \(d\) 的单位。任一齐次元 \(c\in C_e\) 满足 \(c^df^{-e}\in C_0\)。若此元素非零，则 \(c\) 可逆；若为零，则 \(c\) 幂零。因此所有齐次非单位均幂零，它们生成的齐次理想 \(J\) 是幂零元所成的理想。商 \(C/J\) 中每个非零齐次元可逆，故它是整环：两个非零元素的最高非零齐次分量之积非零。因此 \(J\) 是素理想，也是唯一的齐次素理想。沿局部化及剩余域映射收缩 \(J\)，得到唯一避开 \(f\)、且在 \(B\) 中收缩为 \(\mathfrak q\) 的齐次素理想。

这建立了双射。若 \(g\in S_e\)，则在此对应下，\(g\notin\mathfrak p\) 等价于 \(g^d/f^e\notin\mathfrak q\)。于是基本开集 \(D_+(fg)\) 对应 \(B\) 的基本开集 \(D(g^d/f^e)\)，得到同胚。

在该交集上，局部化环满足
\[
(S_{(f)})_{g^d/f^e}\cong (S_{fg})_0.
\]
确实，在 \(S_f\) 中把 \(g^d/f^e\) 变成单位等价于把 \(g\) 变成单位；被倒置的元素次数为零，取零次部分与这次局部化交换。同样从 \(g\) 一侧得到相同环。三重交集上的映射均为分式的限制，所以满足黏合条件。上一节的概形黏合构造给出结构层，且在所有仿射块上的指定限制决定其唯一性。
\end{proof}

\begin{definition}\label{b3:def:projective-space}
设 \(A\) 为交换环，\(r\ge0\)，多项式环 \(A[x_0,\ldots,x_r]\) 中各变量次数为一。概形
\[
\mathbb P_A^r=\operatorname{Proj}A[x_0,\ldots,x_r]
\]
称为 \(A\) 上的 \(r\) 维\term[sheying-kongjian]{射影空间}[projective space]。
\end{definition}

它由 \(r+1\) 个仿射开集覆盖：
\[
D_+(x_i)=\operatorname{Spec}
A[x_0/x_i,\ldots,\widehat{x_i/x_i},\ldots,x_r/x_i]
\cong\mathbb A_A^r.
\]
这些分式代数独立，因为齐次化把它们的多项式关系还原为原变量的关系。对 \(\mathbb P_k^1\)，两张图为 \(k[t]\)、\(k[u]\)，交集为 \(k[t,t^{-1}]\)，黏合规则 \(u=t^{-1}\)。这与双原点直线的恒等黏合不同：射影直线增加的是另一端的点，而非再复制原点。

\ProseParagraphBreak
若 \(I\subseteq S\) 为齐次理想，则 \(\operatorname{Proj}(S/I)\) 在每张图上由
\((S_f)_0\rightarrow((S/I)_f)_0\) 给出商环，因此是 \(\operatorname{Proj}S\) 的闭子概形。齐次方程由此不仅规定闭点的零集，也规定其上的商结构。Proj 的一般构造可对照\cite[Chapter II, Proposition 2.5, pp. 77--78]{Hartshorne1977}。

\subsection{分次模与射影几何}
\label{b3:subsec:2403:03}

设 \(M\) 为分次 \(S\) 模。在 \(D_+(f)\) 上把它取为 \(S_{(f)}\) 模 \((M_f)_0\) 的伴随层；上一定理中的限制同构对模也成立，所以这些层黏合为 \(\operatorname{Proj}S\) 上的拟凝聚层，仍记 \(\widetilde M\)。当同时讨论 Spec 与 Proj 时，必须注明所在空间：前者使用全部 \(M_f\)，后者仅用其次数零部分。

\begin{definition}\label{b3:def:twisting-sheaf}
设 \(S\) 为交换正分次环，\(X=\operatorname{Proj}S\)，\(n\in\mathbb Z\)。沿用次数平移 \(M(n)_j=M_{n+j}\)，定义
\[
\mathcal O_X(n)=\widetilde{S(n)}.
\]
\(\mathcal O_X(1)\) 称为 Serre 的\term[niuqu-ceng]{扭曲层}[twisting sheaf]；对模层 \(\mathcal F\)，记 \(\mathcal F(n)=\mathcal F\otimes_{\mathcal O_X}\mathcal O_X(n)\)。
\end{definition}

\begin{proposition}\label{b3:prop:standard-twists-invertible}
设 \(S\) 为交换正分次环，且由 \(S_1\) 作为 \(S_0\) 代数生成，\(X=\operatorname{Proj}S\)。则各 \(\mathcal O_X(n)\) 都为可逆层；对分次 \(S\) 模 \(M\)，有
\[
\widetilde{M(n)}\cong\widetilde M(n),\qquad
\mathcal O_X(a)\otimes\mathcal O_X(b)\cong\mathcal O_X(a+b).
\]
\end{proposition}
\begin{proof}
次数一元素的 \(D_+(f)\) 覆盖 \(X\)。在这种图上，\(S(n)_{(f)}=(S_f)_n\)，乘以可逆元 \(f^n\) 给出 \(S_{(f)}\cong(S_f)_n\)，故局部秩为一。相同映射给
\((M_f)_0\otimes_{S_{(f)}}(S_f)_n\cong(M_f)_n\)。
这些同构在交集上都由乘法给出，所以黏合。取 \(M=S(a)\) 得最后一式。
\end{proof}

在 \(\mathbb P_k^1\) 上，令 \(t=x_1/x_0\)，两张图中的生成元 \(e_0=x_0^n\)、\(e_1=x_1^n\) 满足 \(e_1=t^ne_0\)。全局截面因而对应
\[
a(t)=t^nb(t^{-1}),\qquad a\in k[t],\quad b\in k[t^{-1}].
\]
当 \(n\ge0\) 时，这正是次数不超过 \(n\) 的多项式，维数为 \(n+1\)；当 \(n<0\) 时只有零截面。特别地，\(\mathcal O(-1)\) 在每点局部自由，却没有非零全局截面，故不可能同构于结构层。

\ProseParagraphBreak
与仿射情形不同，\(M\mapsto\widetilde M\) 在 Proj 上不能由不带次数的全局截面反演。例如 \(S/S_+\) 在每个 \(D_+(f)\) 上局部化为零，故其伴随层为零，尽管这个分次模本身不为零。分次有限模的高次部分、全部扭曲截面与局部上同调之间有更细的关系；在使用这类对应前必须补入相应的有限性与分次假设。扭曲层的局部计算见\cite[Chapter II, Propositions 5.11--5.12, pp. 116--117]{Hartshorne1977}。
